\scanpage{87}

\uline{norme de Schatten} $N(x_1,\ldots x_n)$ satisfait \deleted{à la condition envisagée}\inserted{à l'hypothèse} dans l'énoncé du \deleted{corollaire 1}\inserted{th. 1, 2°} \deleted{\unsure{p.p. la condition de la fonction|application suivante}}\inserted{pour la fonction $f$}. On en conclut, si $(\rho_i)$ \deleted{$\unsure{\rho_i|\rho(x)}$} et $(\sigma_i)$ sont comme dans le \unsure{th.}, :

% slanted
\margin{\textcircled{\unsure{11|13}} $N(\rho_1,\ldots\rho_n)\leq N(\sigma_1,\ldots,\sigma_n)$ et plus généralement, compte tenu d'une remarque qui précède, et en \deleted{\unsure{pour que}} posant pour $1\leq k\leq n$ : $N(\rho_1,\ldots\rho_k)=N(\rho_1,\ldots\rho_k,0,\ldots,0)$ \textcircled{12} $N(\rho_1,\ldots\rho_k)\leq N(\sigma_1,\ldots,\sigma_k)$ pour $1\leq k\leq n$}

\deleted{\unsure{De même}, si on sait que la fonction $(e^{t_i})^p$ étant une \unsure{fonction convexe}} \deleted{\uline{la fonction} $f(r_1,\ldots r_n)=r_1^p+\cdots+r_n^p$, $p$ est un nb réel \uline{quelconque}, satisfait aux \deleted{\unsure{conditions}} conditions \deleted{\unsure{du lemme}} envisagées dans le \uline{corollaire 1}, \deleted{\unsure{donc} dans l'hypothèse b, et même \deleted{\unsure{à toutes les}} conditions envisagées dans l'hypothèse a.}}

Par suite, \deleted{\unsure{on a} \illegible{1} \unsure{inégalités}}\inserted{sous les conditions} de l'hypothèse b, \deleted{\unsure{on}} $(\prod_1^{\unsure{M|N|m}}\rho_i\leq\prod_1^M\sigma_i\ \inserted{(\text{pour }1\leq M\leq n)},\ \prod_1^n\rho_i=\prod_1^n\sigma_i\ne0)$ on aura

\[
\sum_1^n\rho_i^p\leq\sum_1^n\sigma_i^p\qquad(p\text{ réel.})\tag*{(15)}
\]

\deleted{\unsure{et sous les cond}} \inserted{et si $p\geq0$,} on aura, pourvu que $\prod_1^M\rho_i\leq\prod_1^M\sigma_i$ pour \inserted{$1\leq M\leq n$}, \deleted{on aura}

\[
\sum_1^k\rho_i^p\leq\sum_1^k\sigma_i^p\qquad(p\geq0,\ \deleted{\unsure{et}} k=1,\ldots n.)\tag*{\deleted{\unsure{(12)}}\inserted{(16)}}
\]

\deleted{Remarquons d'ailleurs que de ce que $(\rho_i),(\sigma_i)$, \deleted{\unsure{positifs}}\inserted{positifs} \unsure{donnés}, \deleted{\unsure{dès}} satisfaisant aux inégalités (16) pour \deleted{tout} $k$, et \deleted{tout} $p$ entier $\geq0$, on peut \deleted{\unsure{en conclure}}\inserted{réciproquement}}

\deleted{\unsure{Nous pouvons arriver à}} \uline{Th 2} (\uline{Premier Théorème de convexité dans les espaces de Hilbert})

\deleted{\unsure{1)}} Soit $u$ un opérateur compact dans un espace de Hilbert, soit \deleted{\unsure{$\varphi$ la suite}} $f$\inserted{$(r_1,\ldots r_n)$} une fonction \inserted{continue et symétrique} de $n$ arguments $r_i\geq0$, croissante par r. à chaque $r_i$, et telle que $f(e^{t_1},\ldots,e^{t_n})$ soit fonction convexe de $(t_i)\in R^n$. Alors on a (\deleted{\unsure{en prenant} $\rho_i$} (avec les notations indiquées au N° \unsure{?|3}, avant la prop 3)

\[
f(|\lambda_1|,\ldots,|\lambda_n|)\leq f(\rho_1,\ldots\rho_n)\tag*{(17)}
\]

Il suffit en effet de conjuguer le corollaire 1 du lemme 1 \deleted{\unsure{(cas de}} (hypothèse a)) avec le N° \unsure{?|2}, prop 3, corollaire 2. -- En particulier

\uline{Corollaire 1} Si $f\deleted{\unsure{(r)|(\rho)}}$ est une fonction \deleted{\unsure{croissante|convexe}}\inserted{continue} croissante de $r\geq0$ telle que $f(e^t)$ soit convexe, et $u$ un op. compact dans un espace de Hilbert, on a

% archive sheet number: 88
